% ===========================================================================
% 第三章 系统设计与实现
% ===========================================================================
\chapter{系统设计与实现}

第二章给出了\ScaffoldChunGS 的理论约束：全局地图可以随轨迹增长，但GPU只能承载局部工作集；锚点适合作为长期存储单元，子高斯适合作为渲染单元，chunk适合作为空间调度单元。本章在该约束下描述系统实现。章节组织不按模块清单展开，而围绕数据流、生命周期和调度关系展开，核心主线为``块级调度、锚点级存储、子高斯级渲染''。

在工程实现中，系统由C++17、CUDA和LibTorch构成。SLAM前端提供位姿和MappingOperation，Mapper线程将输入帧转化为关键帧和训练任务，GaussianModel维护GPU常驻锚点张量与磁盘chunk文件，AnchorMLP在渲染时把可见锚点展开为子高斯，ScaffoldRenderer执行可微光栅化并将图像损失反传至锚点和MLP。系统扩展能力来自三个解耦：全局地图规模由磁盘chunk承载，GPU显存由active chunk set约束，单帧渲染开销由可见锚点和子高斯数量决定。

\section{总体架构：从帧流到可恢复地图状态}

\subsection{运行时数据流}

\ScaffoldChunGS 的运行时不是一条静态模块链，而是一组围绕地图状态闭环更新的数据流。输入图像首先进入\texttt{SLAMSystemInterface}，前端返回当前帧位姿$\mathbf{T}_{cw}$，并在产生关键帧、局部BA或回环时向MappingOperation队列写入操作。Mapper线程从帧队列和操作队列中取出数据，将关键帧写入\texttt{GaussianScene}，再驱动\texttt{GaussianModel}完成锚点增量、chunk加载、渲染训练和外存淘汰。每次训练迭代只访问由当前关键帧视锥和训练窗口确定的active chunk set，而不是遍历完整地图。

\begin{figure}[htbp]
\centering
\fbox{%
\begin{minipage}{0.94\textwidth}
\centering
\vspace{4pt}
\textbf{Frame Stream}\\
$\downarrow$ \\
\texttt{SLAMSystemInterface}: pose + MappingOperation\\
$\downarrow$ \\
\texttt{ScaffoldMapper}: frame queue + operation dispatch\\
$\downarrow$ \\
\texttt{GaussianScene}: keyframe set \qquad
\texttt{GaussianModel}: anchor tensors + chunk metadata\\
$\downarrow$ \\
active chunk query $\rightarrow$ async load/evict $\rightarrow$ GPU resident anchors\\
$\downarrow$ \\
anchor culling $\rightarrow$ \texttt{AnchorMLP} expansion $\rightarrow$ child Gaussians\\
$\downarrow$ \\
\texttt{ScaffoldRenderer}: differentiable rasterization\\
$\downarrow$ \\
loss backward $\rightarrow$ Adam state update $\rightarrow$ anchor grow/prune $\rightarrow$ chunk persistence
\vspace{4pt}
\end{minipage}}
\caption{系统运行时数据流。图中每个节点对应真实代码对象，箭头表示数据对象的生命周期变化，而不是单纯函数调用顺序。}
\label{fig:ch3_runtime_flow}
\end{figure}

图\ref{fig:ch3_runtime_flow}中的关键约束是：渲染器从不直接面向全局地图工作。它接收的是经过chunk调度和anchor级视锥剔除后的可见锚点切片；子高斯只在一次渲染调用中临时存在；优化器状态与锚点张量绑定，在chunk换出时一同序列化。这种数据边界使系统在长序列运行时保持显存稳定。

\subsection{线程、数据结构与代码边界}

系统中存在三类并发关系。第一，Mapper启动后创建后台线程，执行两阶段建图循环；外部输入通过\texttt{processFrame}写入帧队列，并由\texttt{mutex\_frame\_}保护。第二，Viewer创建独立显示线程，Mapper通过双缓冲\texttt{ViewerFrameData}提交显示快照，Viewer只读取副本，不直接访问训练中的模型。第三，chunk读写并不是单独常驻I/O线程，而是在\texttt{loadChunks}和\texttt{saveChunks}中使用\texttt{std::async}和\texttt{future}按批并行执行，从而缩短训练线程等待磁盘的时间。

\begin{table}[htbp]
\centering
\caption{运行时线程、核心数据结构与代码位置}
\label{tab:runtime_threads}
\begin{tabular}{p{3.0cm}p{5.2cm}p{4.8cm}}
\toprule
\textbf{执行上下文} & \textbf{主要数据结构} & \textbf{代码位置} \\
\midrule
Mapper线程 & 帧队列、MapperState、GaussianScene、GaussianModel & \texttt{mapper\_core.cpp}，\texttt{gaussian\_mapper.h} \\
SLAM操作队列 & MappingOperation、KFTuple、操作队列、队列互斥锁 & \texttt{slam\_system.h}，\texttt{builtin\_adapter.cpp} \\
Chunk并行I/O & ChunkData、磁盘chunk集合、GPU resident chunk集合、future批任务 & \texttt{model\_memory.cpp}，\texttt{model\_storage.cpp} \\
GPU训练路径 & anchor张量、Adam状态、AnchorData、ExpandedGaussians、RenderOutput & \texttt{src/model}，\texttt{src/rendering}，\texttt{src/training} \\
Viewer线程 & ViewerFrameData双缓冲、frame mutex、close condition variable & \texttt{gaussian\_viewer.cpp} \\
\bottomrule
\end{tabular}
\end{table}

表\ref{tab:runtime_threads}说明，本章后续每个系统过程都可以落到真实代码结构：数据进入Mapper线程，形成关键帧和操作；地图状态在GaussianModel内以锚点张量和chunk元数据存在；训练过程生成子高斯并执行反传；Viewer线程只消费快照。这样的线程边界避免了渲染显示、SLAM输入和地图训练之间的直接写冲突。

\subsection{图示设计建议}

本节建议绘制一张黑白系统架构图。图中上层为执行上下文：Input Thread、Mapper Thread、Async I/O Tasks、GPU Kernels、Viewer Thread；下层为数据对象：Frame、MappingOperation、GaussianKeyframe、Anchor Tensor、ChunkData、ExpandedGaussians、RenderOutput。箭头应标注数据类型，例如``pose + image''、``active chunk ids''、``visible anchor slice''和``Adam state''。该图用于替代普通模块框图，重点表达数据流和线程边界。

\section{锚点级地图状态：张量布局与增量生命周期}

\subsection{Anchor tensor layout}

锚点是系统的长期地图单元。一个锚点不是单个点，而是\texttt{GaussianModel}中多组行对齐张量的同一行。设当前GPU常驻锚点数为$N$，每个锚点展开$K$个子高斯，特征维度为$D$，则主要张量布局如表\ref{tab:anchor_tensor_layout}所示。所有字段都以行号作为局部索引，\texttt{anchor\_ids\_}提供跨chunk换入换出的全局身份，\texttt{anchor\_chunk\_ids\_}提供chunk调度所需的空间归属。

\begin{table}[htbp]
\centering
\caption{GPU常驻锚点张量布局}
\label{tab:anchor_tensor_layout}
\begin{tabular}{p{3.5cm}p{2.6cm}p{6.4cm}}
\toprule
\textbf{字段} & \textbf{形状} & \textbf{生命周期含义} \\
\midrule
\texttt{anchor\_} & $[N,3]$ & 锚点中心位置，参与投影、chunk编码和梯度更新 \\
\texttt{offset\_} & $[N,K,3]$ & 子高斯局部偏移，渲染时与基础尺度组合生成子高斯中心 \\
\texttt{anchor\_feat\_} & $[N,D]$ & MLP输入特征，是压缩外观和几何细节的主要载体 \\
\texttt{anchor\_scaling\_} & $[N,6]$ & 前3维控制offset范围，后3维控制子高斯基础尺度 \\
\texttt{anchor\_rotation\_} & $[N,4]$ & 基础四元数旋转 \\
\texttt{anchor\_opacity\_} & $[N,1]$ & 基础不透明度，用于剪枝统计和展开初始化 \\
\texttt{anchor\_chunk\_ids\_} & $[N]$ & 64-bit chunk编码，决定该行属于哪个外存块 \\
\texttt{anchor\_ids\_} & $[N]$ & 全局锚点ID，保证换入换出后的身份连续性 \\
\bottomrule
\end{tabular}
\end{table}

这种布局服务两个工程需求。其一，chunk换出时可以用布尔mask在所有张量上同步切片，形成\texttt{ChunkData}；其二，渲染时可以用可见mask一次性得到\texttt{AnchorData}，再送入AnchorMLP展开函数。因此，张量行对齐是锚点压缩与chunk调度能够耦合的基础。

\subsection{锚点从创建到剪枝的生命周期}

锚点生命周期由初始化、增量添加、稠密化和chunk I/O共同维护。初始建图阶段，Mapper从关键帧深度生成点云，模型按体素大小将点云聚合为锚点，并初始化特征、offset、尺度、旋转和不透明度。增量建图阶段，新关键帧产生的新点云通过\texttt{addAnchors}进入模型，系统进行体素重复过滤，为新增锚点分配递增全局ID，并通过\texttt{computeChunkIds}写入chunk归属。

训练过程中，梯度累积张量、不透明度累积张量和可见计数记录锚点贡献。周期性稠密化由\texttt{adjustAnchors}触发：高梯度区域执行\texttt{anchorGrowing}生成新锚点，低不透明度区域执行\texttt{anchorPruning}移除弱贡献锚点。新增或删除锚点会改变张量行数，因此优化器状态必须同步扩展、剪裁或重建。代码中\texttt{catAnchorsToOptimizer}和\texttt{trainingSetup}负责将新增锚点纳入Adam参数组，并为新增行初始化一阶、二阶矩。

\begin{figure}[htbp]
\centering
\fbox{%
\begin{minipage}{0.92\textwidth}
\centering
\vspace{4pt}
Point/Depth Seeds $\rightarrow$ voxel aggregation $\rightarrow$ Anchor Rows\\
$\downarrow$ \\
assign \texttt{anchor\_ids\_} + \texttt{anchor\_chunk\_ids\_}\\
$\downarrow$ \\
visible training $\rightarrow$ gradient/opacity statistics\\
$\downarrow$ \\
\texttt{anchorGrowing} for high-gradient regions \quad
\texttt{anchorPruning} for low-opacity rows\\
$\downarrow$ \\
optimizer state extension/pruning $\rightarrow$ chunk persistence
\vspace{4pt}
\end{minipage}}
\caption{锚点生命周期。锚点行的创建、增长、剪枝和优化器状态维护必须同步，否则chunk换入换出后无法恢复训练状态。}
\label{fig:anchor_lifecycle}
\end{figure}

\subsection{从锚点存储到子高斯生成}

渲染阶段不会直接把\texttt{GaussianModel}中的全部锚点送入光栅化器。流程为：\texttt{cullVisibleAnchors}先返回可见锚点mask，\texttt{getVisibleAnchors}将行对齐张量切片为\texttt{AnchorData}，随后\texttt{AnchorMLP::expandAnchorsToGaussians}根据锚点特征、观察方向、距离和可选外观嵌入输出\texttt{ExpandedGaussians}。其输出包含\texttt{xyz}、\texttt{color}、\texttt{opacity}、\texttt{scaling}、\texttt{rotation}和\texttt{child\_mask}。其中\texttt{child\_mask}会过滤低不透明度子高斯，使实际进入渲染器的图元数小于$K|\mathcal{A}_{\mathrm{vis}}|$。

该设计使长期地图状态和渲染图元状态分离：chunk文件保存的是锚点行与优化器状态；子高斯是当前视角、当前迭代、当前可见mask下的临时计算结果。反向传播时，梯度从渲染损失经子高斯属性回到MLP权重和可见锚点张量，而不是写回一个全局子高斯表。

\subsection{图示设计建议}

本节建议绘制一张anchor tensor layout图和一张anchor生命周期图。layout图横向画多条等长张量行，使用同一行号$i$连接\texttt{anchor\_}、\texttt{offset\_}、\texttt{anchor\_feat\_}、\texttt{anchor\_chunk\_ids\_}和Adam状态。生命周期图使用状态节点``Seeded''、``Resident''、``Visible/Optimized''、``Grown''、``Pruned''、``Serialized''，箭头标注具体函数名。

\section{块级外存调度：active set、异步I/O与优化连续性}

\subsection{Chunk ID、active chunk与GPU resident set}

Chunk是系统的空间调度单元。锚点位置先映射到\texttt{ChunkCoord}，再编码为64-bit整数。设GPU中当前张量包含的chunk集合为$\mathcal{C}_{\mathrm{res}}$，磁盘中已持久化的chunk集合为$\mathcal{C}_{\mathrm{disk}}$，当前训练需要的active集合为$\mathcal{C}_{\mathrm{act}}$，则系统每次训练前执行的核心关系为
\begin{equation}
\mathcal{C}_{\mathrm{miss}}
=
\mathcal{C}_{\mathrm{act}}
\cap
\mathcal{C}_{\mathrm{disk}}
\setminus
\mathcal{C}_{\mathrm{res}}.
\end{equation}
式中，$\mathcal{C}_{\mathrm{miss}}$表示本次训练前需要从磁盘换入的chunk集合；$\mathcal{C}_{\mathrm{act}}$表示当前帧、训练窗口或回环操作需要访问的active chunk集合；$\mathcal{C}_{\mathrm{disk}}$表示已经持久化到磁盘的chunk集合；$\mathcal{C}_{\mathrm{res}}$表示当前已经常驻GPU张量中的chunk集合；$\cap$表示集合交集，$\setminus$表示集合差。
\texttt{loadChunks}只加载$\mathcal{C}_{\mathrm{miss}}$，并在加载前估计即将进入GPU的锚点数。如果加载后可能超过GPU常驻上限，系统先调用\texttt{evictExcessChunks}，从非保护chunk中选择LRU候选换出。

GPU resident set不是全部地图，而是当前张量中实际保留的锚点行集合：
\begin{equation}
\mathcal{A}_{\mathrm{res}}
=
\{a_i \mid c(a_i)\in\mathcal{C}_{\mathrm{res}}\}.
\end{equation}
式中，$\mathcal{A}_{\mathrm{res}}$表示GPU当前常驻锚点集合；$a_i$表示第$i$个锚点；$c(a_i)$表示锚点$a_i$所属的chunk ID；$\mathcal{C}_{\mathrm{res}}$表示GPU resident chunk集合。该式说明，GPU中的锚点由resident chunk反查得到，而不是由全局地图直接决定。
渲染使用的可见集合进一步满足$\mathcal{A}_{\mathrm{vis}}\subseteq\mathcal{A}_{\mathrm{res}}$。因此，GPU常驻上限约束的是训练态张量规模，\texttt{cullVisibleAnchors}约束的是单帧展开和光栅化规模。

\subsection{Active chunk状态机}

系统中chunk状态由三组信息共同决定：是否在磁盘集合中，是否在GPU resident set或当前新建张量中，是否被当前帧、训练窗口或回环保护。其生命周期可抽象为图\ref{fig:chunk_state_machine}。

\begin{figure}[htbp]
\centering
\fbox{%
\begin{minipage}{0.90\textwidth}
\centering
\vspace{4pt}
New/Resident
$\xrightarrow{\texttt{saveChunks}}$
OnDisk
$\xrightarrow{\texttt{loadChunks}}$
Loaded/Resident
$\xrightarrow{\text{visible or loop}}$
Protected/Active\\[2pt]
Protected/Active
$\xrightarrow{\text{access update}}$
Loaded/Resident
$\xrightarrow{\texttt{findLRUChunks}}$
EvictionCandidate
$\xrightarrow{\texttt{saveAndEvictChunks}}$
OnDisk
\vspace{4pt}
\end{minipage}}
\caption{Active chunk状态机。只有非保护chunk可以进入淘汰候选；淘汰前先序列化锚点参数、统计量和Adam状态。}
\label{fig:chunk_state_machine}
\end{figure}

该状态机解释了为什么chunk不会破坏优化连续性。换出不是丢弃地图，而是将\texttt{ChunkData}写入\schun 文件；换入不是重新初始化，而是把锚点张量、统计张量和Adam状态拼接回GPU参数组。只要chunk文件保存了参数与优化状态，局部区域重新进入active set时就能延续之前的训练轨迹。

\subsection{异步I/O流水线}

\texttt{loadChunks}和\texttt{saveChunks}采用批量并行I/O。读取时，系统把待加载chunk ID分批交给\texttt{std::async}任务，每个任务解析一个\schun 文件，主线程通过\texttt{future}收集\texttt{ChunkData}，最后统一拼接到GPU张量。写出时，\texttt{saveChunks}先在当前线程根据chunk mask提取\texttt{ChunkData}，再并行写入磁盘。

这里的并行粒度是chunk文件，而不是张量元素。原因是单个chunk内部字段较多，包括锚点参数、统计张量和多组Adam矩，逐字段异步会增加同步复杂度；按chunk提交任务可以保持文件原子性和调度语义。代码按硬件并发数限制任务批量，以避免在嵌入式平台上因I/O任务过多挤占训练线程。

\subsection{\schun 文件与Adam状态恢复}

\schun 是chunk的可恢复训练状态文件。每个文件由chunk ID、锚点参数、统计张量和Adam状态组成。表\ref{tab:schun_format}列出主要字段。

\begin{table}[htbp]
\centering
\caption{\schun 文件中的可恢复训练状态}
\label{tab:schun_format}
\begin{tabular}{p{3.3cm}p{8.8cm}}
\toprule
\textbf{字段组} & \textbf{内容} \\
\midrule
锚点参数 & \texttt{anchor}、\texttt{offset}、\texttt{anchor\_feat}、\texttt{anchor\_scaling}、\texttt{anchor\_rotation}、\texttt{anchor\_opacity} \\
身份与归属 & \texttt{exist\_since}、\texttt{anchor\_chunk\_ids}、\texttt{anchor\_ids}、\texttt{chunk\_id} \\
训练统计 & \texttt{offset\_gradient\_accum}、\texttt{offset\_denom}、\texttt{opacity\_accum}、\texttt{anchor\_denom} \\
优化器状态 & 每个锚点参数组的\texttt{exp\_avg}、\texttt{exp\_avg\_sq}和\texttt{step\_counts} \\
\bottomrule
\end{tabular}
\end{table}

换出过程在\texttt{saveAndEvictChunks}中完成。系统先保存loaded和new chunk，再从所有锚点张量中删除对应行，同时裁剪Adam参数状态。换入过程通过append流程完成：系统将chunk数据追加到现有张量末尾，并恢复相应参数组的Adam一阶矩、二阶矩和步数。该流程保证了外存调度只改变数据驻留位置，不改变优化问题本身。

\subsection{图示设计建议}

本节建议绘制三幅黑白图。第一幅为chunk状态机，对应图\ref{fig:chunk_state_machine}；第二幅为异步I/O流水线，显示``chunk id batch''到多个\texttt{future}任务再到\texttt{appendLoadedChunks}的汇聚；第三幅为\schun 文件结构图，将文件分为Header、Anchor Tensors、Statistics和Adam States四段。图中应突出CPU侧调度、磁盘文件和GPU resident set三层关系。

\section{子高斯级渲染与训练闭环}

\subsection{单次训练迭代的数据路径}

训练迭代在Mapper和Trainer中采用一致的数据路径：关键帧采样、视锥剔除、chunk管理、锚点切片、MLP展开、光栅化、损失计算、反向传播、Adam更新和周期性稠密化。该路径如图\ref{fig:training_pipeline}所示。

\begin{figure}[htbp]
\centering
\fbox{%
\begin{minipage}{0.94\textwidth}
\centering
\vspace{4pt}
\texttt{KeyframeSelection}
$\rightarrow$
camera matrices
$\rightarrow$
\texttt{cullVisibleAnchors}\\
$\rightarrow$
active chunk load/evict
$\rightarrow$
\texttt{getVisibleAnchors}
$\rightarrow$
\texttt{AnchorMLP::expandAnchorsToGaussians}\\
$\rightarrow$
\texttt{ScaffoldRenderer::render}
$\rightarrow$
\texttt{computeAllLosses}
$\rightarrow$
\texttt{loss.backward}\\
$\rightarrow$
\texttt{optimizerStep}
$\rightarrow$
\texttt{adjustAnchors}
$\rightarrow$
\texttt{checkMemoryPressure}
\vspace{4pt}
\end{minipage}}
\caption{训练与渲染闭环。GPU只处理当前resident set中的可见锚点，子高斯是MLP展开的临时渲染图元。}
\label{fig:training_pipeline}
\end{figure}

这个闭环中有两个重要资源边界。第一，\texttt{cullVisibleAnchors}在渲染前执行chunk级和anchor级筛选，使MLP前向只作用于$\mathcal{A}_{\mathrm{vis}}$。第二，\texttt{optimizerStep}使用可见mask执行稀疏更新语义，避免不可见锚点在该迭代中被无意义梯度扰动。二者共同保证训练计算集中在当前监督视角相关的地图局部。

\subsection{两级视锥剔除与renderer interaction}

视锥剔除的第一级是chunk级AABB测试。模型根据相机中心和WVP矩阵查询可见chunk ID，必要时调用\texttt{loadChunks}把磁盘chunk换入GPU。第二级是anchor级测试：系统先得到候选锚点mask，再结合投影范围和深度约束得到最终可见mask。该mask一方面保护当前可见chunk免于LRU淘汰，另一方面决定传入渲染器的锚点切片。

渲染器交互由\texttt{ScaffoldRenderer::render}封装。输入包括\texttt{GaussianModel}、可见mask、相机矩阵、图像尺寸和背景色。函数内部获取\texttt{AnchorData}，调用MLP生成\texttt{ExpandedGaussians}，再使用CUDA后端或LibTorch回退后端输出\texttt{RenderOutput}。输出字段包括RGB、depth、alpha和radii，其中radii用于可见性统计和稠密化判断。由于MLP展开和光栅化都在LibTorch计算图内保持可导，损失可以沿\texttt{RenderOutput}反传到锚点特征、offset、尺度、不透明度和MLP权重。

\subsection{损失、反向传播与优化状态}

训练目标由RGB重建、SSIM、深度和各向同性正则组成：
\begin{equation}
\mathcal{L}
=
(1-\lambda_s)\mathcal{L}_{1}
+
\lambda_s\mathcal{L}_{\mathrm{SSIM}}
+
\lambda_d\mathcal{L}_{\mathrm{depth}}
+
\lambda_i\mathcal{L}_{\mathrm{iso}}.
\end{equation}
式中，$\mathcal{L}$为单次训练迭代的总损失；$\mathcal{L}_{1}$为RGB颜色的L1重建损失；$\mathcal{L}_{\mathrm{SSIM}}$为结构相似性损失；$\mathcal{L}_{\mathrm{depth}}$为深度监督损失；$\mathcal{L}_{\mathrm{iso}}$为子高斯尺度各向同性正则；$\lambda_s$、$\lambda_d$和$\lambda_i$分别表示SSIM、深度和尺度正则的权重。
这些损失在\texttt{src/training/losses.cpp}中计算，训练循环调用\texttt{loss.backward()}生成梯度。随后\texttt{GaussianModel::optimizerStep}根据可见mask更新参数组。参数组包括锚点位置、offset、特征、不透明度、尺度、旋转、MLP参数和外观嵌入，对应\texttt{kNumAnchorParamGroups}中的八组状态。由于Adam状态按参数组保存到chunk文件，chunk换出后再换入不会丢失局部动量信息。

锚点生长会改变优化器状态维度。代码中的实现策略是保存已有锚点行的Adam状态，为新增行补零，再重建参数组。这种方式比直接重启优化器更稳定，因为已有锚点的动量和步数得以保留；新增锚点则从零动量开始参与优化，避免污染旧锚点的收敛轨迹。

\subsection{显存稳定性的来源}

单次训练迭代的GPU显存主项可写为
\begin{equation}
M_{\mathrm{iter}}
\approx
B_A|\mathcal{A}_{\mathrm{res}}|
+
M_{\mathrm{MLP}}
+
M_{\mathrm{render}}(K|\mathcal{A}_{\mathrm{vis}}|)
+
M_{\mathrm{frame}}.
\end{equation}
式中，$M_{\mathrm{iter}}$表示单次训练迭代的主要GPU显存占用；$B_A$表示单个锚点及其训练态状态的平均字节数；$\mathcal{A}_{\mathrm{res}}$表示GPU resident锚点集合；$M_{\mathrm{MLP}}$表示共享MLP参数及其优化状态占用；$M_{\mathrm{render}}(\cdot)$表示渲染中间缓存随输入子高斯数量变化的显存开销；$K$为每个锚点展开的子高斯数量；$\mathcal{A}_{\mathrm{vis}}$表示当前帧可见锚点集合；$M_{\mathrm{frame}}$表示输入图像、深度图和监督缓存占用。
系统通过\texttt{max\_anchors\_in\_memory\_}约束$|\mathcal{A}_{\mathrm{res}}|$，通过两级视锥剔除约束$|\mathcal{A}_{\mathrm{vis}}|$，通过chunk外存让全局地图规模$|\mathcal{A}|$转移到磁盘。因此，即使全局地图继续增长，GPU训练态显存仍围绕active working set波动，而不是随轨迹长度线性增长。

\subsection{图示设计建议}

本节建议绘制一张pipeline图和一张反向传播图。pipeline图按图\ref{fig:training_pipeline}组织，每个节点标注函数名。反向传播图从\texttt{RenderOutput}回指\texttt{ExpandedGaussians}、\texttt{AnchorMLP}和anchor tensor rows，用虚线表示梯度流，用实线表示前向数据流。图中应明确标注child Gaussian为runtime object，anchor tensor为persistent trainable state。

\section{SLAM驱动的Mapper pipeline与同步机制}

\subsection{两阶段Mapper生命周期}

\texttt{ScaffoldMapper}是系统顶层编排器，其生命周期由\texttt{MapperState}表示，包括\texttt{kIdle}、\texttt{kInitializing}、\texttt{kTracking}、\texttt{kLoopClosing}和\texttt{kShuttingDown}。启动后，\texttt{start}创建Mapper线程并调用\texttt{run}。系统先进入Phase 1初始建图：累积若干关键帧，估计或读取深度，调用\texttt{sampleGaussians}和\texttt{seedGaussians}生成初始锚点，再执行\texttt{trainingSetup}建立Adam参数组。满足初始条件后进入Phase 2增量映射：持续消费帧队列和MappingOperation队列，执行局部训练、回环重优化和尺度修正。

\begin{figure}[htbp]
\centering
\fbox{%
\begin{minipage}{0.94\textwidth}
\centering
\vspace{4pt}
\textbf{Phase 1: Initial Mapping}\\
frame queue $\rightarrow$ track/pose $\rightarrow$ GaussianKeyframe $\rightarrow$
depth seeds $\rightarrow$ initial anchors $\rightarrow$ training setup\\[4pt]
\textbf{Phase 2: Incremental Mapping}\\
MappingOperation queue $\rightarrow$ combine operations $\rightarrow$
LocalBA / LoopClosingBA / ScaleRefinement\\
$\downarrow$\\
active chunk update $\rightarrow$ train iterations $\rightarrow$ viewer snapshot $\rightarrow$ save chunks
\vspace{4pt}
\end{minipage}}
\caption{Mapper两阶段生命周期。Phase 1建立可训练锚点地图，Phase 2由SLAM操作队列持续驱动局部更新和回环修正。}
\label{fig:mapper_lifecycle}
\end{figure}

与普通离线3DGS训练不同，Mapper的输入不是固定训练集，而是不断变化的关键帧流和操作流。每个\texttt{MappingOperation}携带一组\texttt{KFTuple}，其中包含关键帧ID、相机ID、位姿、主图像、回环标志和辅助图像。Mapper将这些元组转化为\texttt{GaussianKeyframe}，写入带mutex保护的\texttt{GaussianScene}，再由关键帧采样器选择训练监督。

\subsection{SLAM前端如何驱动地图更新}

\texttt{SLAMSystemInterface}抽象了三类跟踪入口：\texttt{trackMonocular}、\texttt{trackStereo}和\texttt{trackRGBD}。内置\texttt{BuiltinSLAMAdapter}使用ORB特征、PnP和BoVW回环作为自包含实现；接口也允许替换为更完整的SLAM前端。无论前端实现如何，Mapper只依赖两个输出：当前帧位姿和MappingOperation队列。

LocalMappingBA表示局部窗口需要补充训练。Mapper通过LocalBA批处理函数处理相关关键帧，生成新锚点或更新局部窗口。LoopClosingBA表示位姿图或回环检测已经改变历史区域的一致性约束。Mapper进入\texttt{kLoopClosing}状态，暂停普通帧接收，临时扩大GPU锚点上限，保护回环相关chunk，并执行额外回环优化迭代。ScaleRefinement则处理尺度变化对地图坐标和锚点分布的影响。

\subsection{多线程同步与数据一致性}

系统同步机制围绕``共享状态最小化''设计。输入帧进入\texttt{frame\_queue\_}时由\texttt{mutex\_frame\_}保护；Mapper状态由\texttt{mutex\_state\_}保护；\texttt{GaussianScene}内部的关键帧map由\texttt{mutex\_kfs\_}保护；\texttt{BuiltinSLAMAdapter}的MappingOperation队列由\texttt{mutex\_queue\_}保护；Viewer通过\texttt{frame\_mutex\_}读取双缓冲快照。GPU训练路径本身集中在Mapper线程中执行，避免多个线程同时写\texttt{GaussianModel}张量。

\begin{figure}[htbp]
\centering
\fbox{%
\begin{minipage}{0.94\textwidth}
\centering
\vspace{4pt}
Input Thread $\rightarrow$ \texttt{processFrame} $\rightarrow$ \texttt{frame\_queue\_}\\
\hfill $\downarrow$ Mapper Thread consumes frames and operations\\
SLAM Adapter $\rightarrow$ \texttt{op\_queue\_} $\rightarrow$ \texttt{combineMappingOperations}\\
\hfill $\downarrow$ train/render/update GaussianModel\\
Mapper Thread $\rightarrow$ \texttt{submitToViewer} $\rightarrow$ Viewer double buffer $\rightarrow$ Viewer Thread
\vspace{4pt}
\end{minipage}}
\caption{Mapper、SLAM队列与Viewer之间的同步关系。训练张量只由Mapper线程写入，其他线程通过队列或快照传递数据。}
\label{fig:thread_sync}
\end{figure}

这种同步方式使系统在工程上保持清晰的数据所有权：SLAM前端拥有位姿估计和操作生成，Mapper拥有地图训练状态，I/O任务拥有单个chunk文件的临时读写，Viewer拥有显示快照。chunk换入换出、锚点增删和Adam状态恢复均在Mapper控制路径下完成，避免了异步I/O直接修改训练中张量的风险。

\subsection{Loop closure与active chunk保护}

回环闭合对大场景系统的影响不止是更新相机轨迹。它可能重新激活已经换出到磁盘的历史chunk，并要求这些chunk在若干训练迭代内保持常驻。系统通过两个机制处理该问题。第一，LoopClosingBA期间临时放大内存上限，避免受影响区域刚加载又被LRU淘汰。第二，回环相关关键帧可见的chunk进入保护集合，\texttt{evictExcessChunks}只从保护集合外选择候选。这样，回环优化期间地图一致性修正不会被外存调度打断。

\subsection{图示设计建议}

本节建议绘制一张时序图和一张Mapper状态图。时序图包含Input、SLAMSystemInterface、Operation Queue、ScaffoldMapper、GaussianModel、Renderer和Viewer七条泳道，重点标注LocalBA和LoopClosingBA两种路径。状态图以\texttt{MapperState}为节点，标出初始建图、正常跟踪、回环闭合和关闭之间的转换条件。

\section{实现索引与配置边界}

\subsection{配置如何控制资源约束}

系统资源边界由\texttt{scaffold\_chunks.yaml}控制，并在\texttt{trainer.cpp}中映射到锚点、chunk、优化、相机和关键帧五组配置。关键配置不是简单超参数，而是直接改变数据流粒度：offset数量决定每个锚点生成的子高斯数，特征维度决定MLP输入和压缩容量，chunk边长决定I/O空间粒度，GPU锚点上限决定resident set规模，损失权重决定反向传播对颜色、深度和形状的约束比例。

\begin{table}[htbp]
\centering
\caption{配置项与系统资源边界}
\label{tab:config_behavior}
\begin{tabular}{p{3.6cm}p{2.1cm}p{6.3cm}}
\toprule
\textbf{配置项} & \textbf{默认值} & \textbf{影响的数据流或资源边界} \\
\midrule
\texttt{n\_offsets} & 5 & 控制$K$，影响子高斯展开规模和渲染负载 \\
\texttt{feat\_dim} & 32 & 控制锚点特征容量和MLP输入维度 \\
\texttt{voxel\_size} & 0.01 & 控制锚点初始化密度和增量生长分辨率 \\
\texttt{chunk\_size} & 20.0 & 控制chunk AABB、I/O粒度和active set大小 \\
\texttt{max\_anchors} & 120000 & 控制GPU resident anchor上限和LRU触发阈值 \\
\texttt{lambda\_dssim} & 0.2 & 控制结构相似性损失权重 \\
\texttt{lambda\_depth} & 0.1 & 控制深度监督对地图几何的约束强度 \\
\texttt{lambda\_isotropic} & 0.01 & 控制子高斯尺度正则强度 \\
\bottomrule
\end{tabular}
\end{table}

\subsection{源文件与系统职责}

表\ref{tab:fileorg}给出本章叙事与真实源文件的对应关系。该表的作用是帮助定位实现，而不是替代系统数据流说明。

\begin{table}[htbp]
\centering
\caption{核心源文件与系统职责}
\label{tab:fileorg}
\begin{tabular}{p{4.2cm}p{7.6cm}}
\toprule
\textbf{文件} & \textbf{系统职责} \\
\midrule
Mapper文件族 & Mapper线程、两阶段建图、MappingOperation分发、训练闭环编排 \\
GaussianModel文件族 & 锚点张量、初始化、增量添加、全局ID与chunk归属 \\
Memory/Storage文件族 & active chunk、LRU、并行I/O、\schun 序列化与Adam状态恢复 \\
AnchorMLP文件族 & 可见锚点切片到子高斯的动态展开 \\
Optimization文件族 & Adam参数组、稀疏更新、锚点生长剪枝、优化器状态扩展 \\
FrustumCuller文件族 & chunk级和anchor级视锥剔除 \\
Renderer文件族 & ScaffoldRenderer、CUDA/LibTorch渲染后端、RenderOutput \\
SLAM接口文件族 & SLAM接口、MappingOperation队列、内置跟踪与回环适配 \\
Scene/Keyframe文件族 & 关键帧容器、线程安全访问、训练关键帧采样 \\
Viewer文件族 & Viewer线程、双缓冲快照、OpenGL可视化 \\
\bottomrule
\end{tabular}
\end{table}

\section{本章小结}

本章从系统论文角度重构了\ScaffoldChunGS 的设计与实现。系统以Mapper线程为调度中心，以SLAM位姿和MappingOperation驱动关键帧生命周期，以chunk作为CPU/GPU/磁盘之间的调度单位，以anchor tensor row作为可持久化训练状态，以child Gaussian作为运行时渲染图元。GPU显存稳定来自active chunk set和\texttt{max\_anchors\_in\_memory}约束，渲染效率来自chunk级与anchor级两级剔除，优化连续性来自\schun 文件对锚点参数、统计量和Adam状态的联合持久化。下一章将围绕压缩率、显存占用、I/O吞吐、渲染质量和SLAM运行效率验证这些系统设计是否达到预期效果。
